\section{An observation with randomly displaced windows}\label{sec:layout}

The nested words now have the overlap needed to hide their motif locations.  We construct an auxiliary observation that reports the tree of windows and the side traces while retaining the hidden fair orientation bits.  Its projection to a genuine trace is justified before any conditioning.  The score proof will then condition only on the reported side data.

A related auxiliary experiment in the probability-vector model of \citet[Section 3, Lemmas 6--8]{RivkinValiantValiant2025} exposes a binomially translated middle segment and separately deleted sides.  That model draws fresh Bernoulli bits from the probability vector before each trace.  Here the input is a fixed deterministic word, and the construction below proves the projection needed for that input.

\subsection{The index tree and the outer window}

Use a deterministic tree common to both hypotheses.  Its root $\varnothing$ is the sole level-$b$ vertex.  A level-$j$ vertex has $T_j$ children at level $j-1$, indexed by appending $h\in\{1,\ldots,T_j\}$ to its own index.  Thus an index records positions relative to successive sibling groups.  Level-zero vertices are leaves.  The entire experiment described below is redrawn independently for each observation.

At the root, place a central run of $T_b$ level-$(b-1)$ atoms between $J$ pairs of pure $U_{b-1}^+$ runs.  For each pair use an independent fair bit $E$, giving counts $h_{b-1}+E$ on the left and $h_{b-1}+1-E$ on the right.  Use the paired flag mechanism from Section~\ref{sec:tools}, with hidden probability $\omega^2$, at the overlap proved in Lemma~\ref{lem:nested-overlap}.  Report one unseparated ordinary trace per side run, its flag, and its orientation bit when revealed.  The left runs have fixed pair-number order before the center, and the right runs have the corresponding fixed order afterward.

The root motif at the word label $\lfloor P_b/2\rfloor+a$ has relative label $X+a$ in this central run, where
\begin{equation}\label{eq:outer-anchor}
 X=\lfloor P_b/2\rfloor-Jh_{b-1}-\sum E
       =T_b/2+O(J+1).
\end{equation}
This is the coordinate of a fixed motif relative to the moving central window.  Since $J=o(T_b)$ and $L=o(T_b)$, every $X+a$, $a\in\mathcal I$, lies in $[T_b/3,2T_b/3]$ for all orientations and all sufficiently large instances.

\subsection{One group of child windows}\label{sec:group-partition}

Fix $1\le j<b$ and a level-$(j+1)$ vertex $G$.  Its $T=T_{j+1}$ children $Gu$, $1\le u\le T$, represent a planned consecutive sequence of $T$ level-$j$ units.  We partition that sequence into successive parts, measured in level-$(j-1)$ atoms.  Part $u$ contains a window of $T_j$ atoms for the children of $Gu$, with nominal pre and post counts
\[
 A_j=\lfloor(P_j-T_j)/2\rfloor,
 \qquad P_j-T_j-A_j.
\]
The parts may have different actual lengths, but their total length will remain $TP_j$.

Reserve $t_j$ local pairs in each part.  For $1\le v\le t_j$, a fair bit $E^{\rm loc}_{u,v}$ gives a pre chunk of $h_{j-1}+E^{\rm loc}_{u,v}$ pure $U_{j-1}^+$ atoms and a post chunk of $h_{j-1}+1-E^{\rm loc}_{u,v}$ such atoms.  Also reserve global pairs between distant parts.  For every $1\le a\le T/4$ and $1\le v\le t_j$, a fair bit $E^{\rm glob}_{a,v}$ gives a pre chunk of $h_{j-1}+E^{\rm glob}_{a,v}$ atoms in part $a$ and a post chunk of $h_{j-1}+1-E^{\rm glob}_{a,v}$ atoms in part $T+1-a$.  All these bits are independent.

Treat each reserved chunk as having nominal count $h_{j-1}+1/2$.  Complete every pre and post region to its nominal count by one pure-$U_{j-1}^+$ filler of deterministic count, and order its chunks and filler by fixed relative IDs.  The filler counts are nonnegative: each region reserves at most $2t_j$ chunks, while the choice of $t_j$ reserves at most $P_j/50$ atoms and $T_j=o(P_j)$.  They are integers because each reserved collection has the even number $t_j$ of chunks; a region contains zero, one, or two such collections.  Generate each reserved pair of chunks with the joint flag mechanism at the overlap of $U_{j-1}^+$; generate each filler as an independent ordinary trace.  Each chunk or filler reports one unseparated trace, with the flags and revealed bits for paired chunks.  The surrogate uses these pure-run laws for every orientation, even when the corresponding locations would not be pure in the true word.

The exact prefix count explains the two displacement scales.  Define
\begin{align}
 H_{Gu}&=\sum_{v=1}^{t_j}(E^{\rm loc}_{u,v}-1/2),\nonumber\\
 V_{Gu}&=\sum_{a=1}^{T/4}\sum_{v=1}^{t_j}
 \bigl({\bf1}_{a\le u}-{\bf1}_{T+1-a<u}\bigr)
       (E^{\rm glob}_{a,v}-1/2),\qquad Y_{Gu}=V_{Gu}+H_{Gu}.
 \label{eq:prefix-offset}
\end{align}
Then the number of lower atoms before the window of $Gu$, counted from the start of the group, is
\begin{equation}\label{eq:window-start}
 (u-1)P_j+A_j+Y_{Gu}.
\end{equation}
In \eqref{eq:prefix-offset}, a global pre chunk is included when its part $a$ is at or before $u$, whereas its post chunk is included only when part $T+1-a$ is strictly before $u$.  Completed global pairs cancel.  Every earlier local pair has also been completed; only the current local pre chunks contribute $H_{Gu}$.  This proves \eqref{eq:window-start}.  Each included collection contributes the even number $t_j$ of half-integer deviations, so the displacements are integers.  After the last post region all pairs have cancelled, proving the total count $TP_j$.

For the middle half of the group, every global pre chunk has appeared and no global post chunk has appeared.  Thus
\begin{equation}\label{eq:global-plateau}
 V_{Gu}=V_G:=\sum_{a=1}^{T/4}\sum_{v=1}^{t_j}
       (E^{\rm glob}_{a,v}-1/2)
       \quad\text{when }T/4<u\le3T/4.
\end{equation}
The local sums $H_{Gu}$ are independent across $u$ and independent of the global bits.  The global sums outside this interval are partial sums of the same group's bits.  Distinct groups and depths use independent mechanisms.

\subsection{The surrogate law and its projection}

We now apply this group partition recursively.  The root center is the first group of $T_b$ planned level-$(b-1)$ units.  A window at any later vertex is partitioned into the group for its children by the construction above.

Define type labels for every orientation, starting from $\tau_\varnothing=+1$ or $-1$ under the two hypotheses.  At the root put
\[
 \tau_h=1-\chi(h-X)+\sigma_{h-X}\tau_\varnothing.
\]
For every other level-$j$ vertex $I$, $j\ge1$, put
\begin{equation}\label{eq:surrogate-types}
 \tau_{Ih}=1-\chi(h-w_j+Y_I)+\sigma_{h-w_j+Y_I}\tau_I,
 \qquad w_j=\lfloor P_j/2\rfloor-A_j=T_j/2.
\end{equation}
The last equality uses the evenness of $T_j$.  These rules place the motif labels that fit in each window and use type $+$ at its other labels.  They define a surrogate even on a failed layout, when an intended motif label falls outside its window; all non-leaf pieces still use the pure laws already specified.

Conditional on all bits, sample the leaf outputs independently of one another and of the entire collection of non-leaf flags and emissions, using law $\mu_{i-1}$ for type $+1$ and $\nu_{i-1}$ for type $-1$.  The complete report consists of these separate leaf outputs, all side outputs, flags, and visible bits, with their relative tree and segment IDs.  It contains no hidden type labels, absolute word labels, or hidden boundary positions.  Let $D_0$ denote all of this report except the complete old oracle payload at each leaf; in particular, any metadata inside an old payload are excluded from $D_0$.

Let $\widetilde\mu_i,\widetilde\nu_i$ be the two unconditioned surrogate laws.  A common kernel $\mathcal K_{\rm tree}$ applies $\mathcal K_{i-1}$ independently to every leaf output and concatenates the resulting binary words and side traces in the chronological order specified by the IDs: pre pieces, the recursively ordered window contents, and post pieces in each part, with the outer side runs in their stated order.

Define the fit event, which depends on the orientation bits, by
\begin{equation}\label{eq:fit_event}
 \mathcal F=\bigl\{|Y_I|\le T_j/8
       \text{ for every level-$j$ vertex }I,\ 1\le j<b\bigr\}.
\end{equation}

\begin{lemma}[Simulation of a genuine trace]\label{lem:tree-simulation}
For the unconditioned surrogate laws and the common projection above,
\begin{equation}\label{eq:projection_tree}
\begin{aligned}
 \TV(\mathcal K_{\rm tree}\widetilde\mu_i,\Tr(U_b^+))
 &\le o(\Delta)+\left(\prod_{j=1}^bT_j\right)E_{i-1},\\
 \TV(\mathcal K_{\rm tree}\widetilde\nu_i,\Tr(U_b^-))
 &\le o(\Delta)+\left(\prod_{j=1}^bT_j\right)E_{i-1}.
\end{aligned}
\end{equation}
All error terms are uniform in $i$.
\end{lemma}

\begin{proof}
By \eqref{eq:prefix-offset}, each $Y_I$ is a centered sum of at most $(T_{j+1}+1)t_j$ independent fair bits with coefficients of magnitude at most one.  Hoeffding's inequality \cite[Theorem~2]{Hoeffding1963} and \eqref{eq:kj},\eqref{eq:scale_t} give
\[
 \Prob(|Y_I|>T_j/8)
 \le 2\exp\!\left(-c\frac{T_j^2}{T_{j+1}t_j}\right)
 \le 2\exp(-c'N^\epsilon)
\]
for absolute positive constants at all sufficiently large instances.  Equation~\eqref{eq:power_gap} bounds the number of vertices by $O_B(N)$.  A union bound therefore gives
\begin{equation}\label{eq:fit-probability}
 \Prob(\mathcal F^c)\le\exp(-N^{\epsilon/2})=o(\Delta).
\end{equation}
The last comparison follows uniformly from Lemma~\ref{lem:scale-margins}: $\epsilon\log N\ge\epsilon\log N_0$ dominates $\log\log n+\log(d_D+B+1)$, so $N^{\epsilon/2}$ dominates $\log\Delta^{-1}$.

We verify what a fitting orientation does to a group of true units.  Suppose its incoming sequence is $T$ actual level-$j$ words with their prescribed types.  In the lower-atom coordinates of that sequence, the possible motif labels of unit $u$ form
\[
 \mathcal M_u=(u-1)P_j+\lfloor P_j/2\rfloor+\mathcal I,
\]
and its window occupies the integer interval
\[
 \mathcal W_u=[(u-1)P_j+A_j+Y_{Gu}+1,
                (u-1)P_j+A_j+Y_{Gu}+T_j].
\]
The serialized nonnegative pre/window/post parts make these windows disjoint and ordered.  On $\mathcal F$, every motif label has window coordinate $w_j+a-Y_{Gu}\in[1,T_j]$, since $w_j=T_j/2$, $|Y_{Gu}|\le T_j/8$, and $L=o(T_j)$.  Hence $\mathcal M_u\subseteq\mathcal W_u$ for every $u$.  Disjointness now implies that no window contains a different unit's motif, and the complement of all windows contains no motif label at all.  Thus all chunks and fillers are truly pure, even if a part crosses a nominal unit boundary, and the types inside each window are exactly \eqref{eq:surrogate-types}.

At the root, \eqref{eq:outer-anchor} puts every motif label in the center and gives its correct child type.  Applying the preceding group argument successively down the tree proves agreement of every surrogate type with the true lower atom on $\mathcal F$.  Given such orientation bits, averaging the flags gives independent correct trace laws on all side pieces.  The old kernels at the leaves add at most one error $E_{i-1}$ per leaf.  Couple the unconditioned experiments using these bits and flags; charge cost one on $\mathcal F^c$.  Equation~\eqref{eq:fit-probability} and the leaf count give \eqref{eq:projection_tree}.
\end{proof}

\subsection{The information retained after the report}\label{sec:metadata}

The law of $D_0$ is common to the two surrogate hypotheses, because every non-leaf emission is generated from a pure run and all IDs are fixed.  Conditional on its full value, every unrevealed orientation is still a fair bit.  These hidden bits are independent across pairs, groups, and depths: the report of a hidden pair has the common product law, while a non-hidden pair reveals its bit, and the pair mechanisms were generated independently.  Visible contributions are fixed by the report.  Consequently $X$, each local sum, and each global sum have exactly the affine hidden-bit forms given above.  Global offsets within a group retain their prescribed dependence.

Call a metadata value typical when all of the following hold:
\begin{itemize}
\item the root has at least $\omega^2J/2$ hidden orientations;
\item for every $1\le j<b$, every local collection at a level-$j$ vertex has at least $\omega^2t_j/2$ hidden orientations, and every global collection in a group of $T_{j+1}$ children has at least $\omega^2(T_{j+1}/4)t_j/2$;
\item $\Prob(\mathcal F^c\mid D_0)\le\exp(-N^{\epsilon/4})$.
\end{itemize}
These conditions fail with probability $o(\Delta)$.  For the hidden counts, apply multiplicative binomial concentration to the independent flags and then a union bound.  Every required expectation is at least $e^{-O_B((D+1)\Lambda)}N^\epsilon$, and Lemma~\ref{lem:scale-margins} makes the resulting tail negligible relative to $\Delta$.  For the last condition, apply Markov's inequality to \eqref{eq:fit-probability}.

The score analysis in the next section fixes one exact typical value of $D_0$ and averages over all its remaining hidden bits.  It does not condition those bits on $\mathcal F$.
